\documentclass{article}

\usepackage{algorithmicx}
\usepackage[ruled]{algorithm}
\usepackage{algpseudocode}
\usepackage{algpascal}
\usepackage{algc}

\newcommand{\alg}{\texttt{algorithmicx}}
\newcommand{\old}{\texttt{algorithmic}}
\newcommand{\euk}{Euclid}
\newcommand\ASTART{\bigskip\noindent\begin{minipage}[b]{0.5\linewidth}}
\newcommand\ACONTINUE{\end{minipage}\begin{minipage}[b]{0.5\linewidth}}
\newcommand\AENDSKIP{\end{minipage}\bigskip}
\newcommand\AEND{\end{minipage}}

\title{BookML test: the \texttt{algorithmicx} package}
\author{Sz\'asz J\'anos\\szaszjanos@users.sourceforge.net}

\usepackage{bookml/bookml}
\usepackage{framed}
\AtBeginDocument{\begin{framed}
  \begin{description}
    \item[BookML] \BookMLversion
    \item[\LaTeXML] \LaTeXMLversion{} (\csname bml@generator@engine\endcsname)
    \item[\LaTeX] \fmtversion
  \end{description}
\end{framed}}


\begin{document}

\alglanguage{pseudocode}

\begin{algorithm}[H]
\caption{\euk's algorithm}\label{euclid}
\begin{algorithmic}[1]
\Procedure{\euk}{$a,b$}\Comment{The g.c.d. of a and b}
   \State $r\gets a\bmod b$
   \While{$r\not=0$}\Comment{We have the answer if r is 0}
      \State $a\gets b$
      \State $b\gets r$
      \State $r\gets a\bmod b$
   \EndWhile\label{euclidendwhile}
   \State \Return $b$\Comment{The gcd is b}
\EndProcedure
\end{algorithmic}
\end{algorithm}

\begin{algorithmic}[1]
\Repeat
\Comment{forever}
\State this\Until{you die.}
\end{algorithmic}

\end{document}
